Giving tail-risk models the information of markets that have already closed improves forecasts substantially, and a two-parameter extension of GJR-GARCH captures most of the gain. The deep models we designed for the task did not. These were an attention-based contagion network, a neural GARCH with provable stability, and versions trained with an efficiency-improving composite score. They react to the news, but they do not turn it into better forecasts. The result holds in a development sample and, after a pre-registered freeze, in twelve markets that no model had seen. Combinations of econometric and neural forecasts remain the most accurate and the best calibrated, in line with our companion paper.

The mechanism, as far as our evidence identifies it, is statistical efficiency. FZ0 learns from tail days, while a likelihood learns volatility dynamics from every day. The correct functional form, foreign squared news entering the variance additively, is easy to impose and hard to discover from eight-year windows. The broader lesson for machine learning in finance is that, in low signal-to-noise problems, the inductive bias and the information set deserve as much attention as the architecture. A deep model should be compared with the best structured model \emph{given the same information}, not with a model that has less.

\paragraph{Limitations and extensions.} We use daily close-to-close returns of equity indices and a one-day horizon. Intraday data would sharpen the timing of transmission, and other asset classes would test its generality. Larger pooled samples, for example thousands of individual stocks, may change the efficiency trade-off in favour of flexible models. The finding that tail risk flows mainly from the West to Asia overnight invites a structural follow-up: whether the transmitted component of tail risk is priced, in the spirit of \citet{kelly2014tail} and \citet{bali2009intertemporal}.
